\section{Conclusion}

This project was both challenging and rewarding. I had prior exposure to quantitative finance, but credit derivatives were new to me, so building everything from first principles was a significant learning experience. Working through reduced–form credit modelling, hazard rates, and the HJM framework gave me a much clearer understanding of how credit term structures behave and how they feed into pricing.

A major part of the project involved implementing complex payoff logic in Python. The range–accrual CLN required careful structuring of the simulation code and forced me to think clearly about discounting, survival weighting, recovery modelling, and pathwise state variables such as CDS spreads. Writing all of this from scratch was difficult, but it made the numerical experiments feel much more meaningful.

I also used ChatGPT as a technical assistant throughout the project, particularly for debugging plotting code and refining the appearance of the figures. The IEEE–style funnel plots and curve comparison graphs were produced with the help of ChatGPT’s suggestions on styling, layout, and Matplotlib configuration. This significantly improved the clarity and presentation quality of the results, and I have cited these interactions directly in my code where appropriate.

The simulation results validated the modelling choices. The HJM engine repriced the input curves accurately, the Monte Carlo behaviour matched theoretical expectations, and both variance reduction techniques improved efficiency relative to the base case. Antithetic variates delivered strong and consistent variance reduction. The control variate provided only modest improvement, which reflects the limited sensitivity of the exotic payoff to the vanilla CLN under the chosen barrier levels. Even so, the variance reductions across all methods were clean and visible in the plots.

In terms of future work, there are several natural extensions. Importance sampling would be a valuable technique to explore, especially in credit settings where tail behaviour and rare events matter. Quasi–Monte Carlo methods would also be interesting to test, particularly with JAX acceleration. Another important next step would be to reimplement the entire simulation engine in JAX. This would allow vectorised path simulation, automatic differentiation, and GPU acceleration, which would significantly improve both performance and flexibility.

The project was demanding, especially given the scope for a single person, but it pushed me to learn a great deal about both credit modelling and simulation methodology. Despite the complexity, I enjoyed working through the material and seeing the numerical results come together in a coherent way.
